\documentclass[12pt,a4paper]{article}

% --- Packages ---
\usepackage[T1]{fontenc}
\usepackage[utf8]{inputenc}
\usepackage[margin=2.5cm]{geometry}
\usepackage{lmodern}
\usepackage{microtype}
\usepackage{graphicx}
\usepackage{float}
\usepackage{booktabs}
\usepackage{tabularx}
\usepackage{array}
\usepackage{multirow}
\usepackage{amsmath}
\usepackage{amssymb}
\usepackage{bm}
\usepackage{xcolor}
\usepackage{listings}
\usepackage{caption}
\usepackage{subcaption}
\usepackage{fancyhdr}
\usepackage{titlesec}
\usepackage{abstract}
\usepackage{parskip}
\usepackage{enumitem}
\usepackage[hidelinks]{hyperref}

% --- Colours ---
\definecolor{myblue}{RGB}{30,90,160}
\definecolor{mygray}{RGB}{245,245,245}
\definecolor{codegray}{rgb}{0.5,0.5,0.5}
\definecolor{codepurple}{rgb}{0.58,0,0.82}

% --- Hyperref setup ---
\hypersetup{
    colorlinks=true,
    linkcolor=myblue,
    citecolor=myblue,
    urlcolor=myblue,
    pdftitle={ECG Arrhythmia Detection using Deep Learning},
    pdfauthor={Revant Sinha},
}

% --- Header / Footer ---
\pagestyle{fancy}
\fancyhf{}
\renewcommand{\headrulewidth}{0.4pt}
\fancyhead[L]{\small ECG Arrhythmia Detection via Deep Learning}
\fancyhead[R]{\small \thepage}
\fancyfoot[C]{\small Biomedical Instrumentation Project Report}

% --- Section formatting ---
\titleformat{\section}{\large\bfseries\color{myblue}}{\thesection}{1em}{}[\titlerule]
\titleformat{\subsection}{\normalsize\bfseries}{\thesubsection}{1em}{}

% --- Code listing style ---
\lstset{
    backgroundcolor=\color{mygray},
    basicstyle=\ttfamily\small,
    breaklines=true,
    captionpos=b,
    commentstyle=\color{codegray},
    keywordstyle=\color{myblue}\bfseries,
    stringstyle=\color{codepurple},
    numbers=left,
    numberstyle=\tiny\color{codegray},
    numbersep=5pt,
    frame=single,
    framerule=0.4pt,
    language=Python,
    showstringspaces=false,
    tabsize=4,
    literate={<}{{\textless}}1 {>}{{\textgreater}}1,
}

% ====================================================================
\begin{document}
% ====================================================================

% --- Title page ---
\begin{titlepage}
    \centering
    \vspace*{2cm}
    {\LARGE\bfseries\color{myblue}
        Deep Learning-Based Arrhythmia Detection\\[0.4em]
        from ECG Signals\\[0.4em]
        \large A Comparative Study of Custom CNN vs.\ Transfer Learning
    \par}
    \vspace{1cm}
    \rule{\textwidth}{1.5pt}
    \vspace{0.5cm}

    {\large\bfseries Biomedical Instrumentation --- Project Report\par}
    \vspace{0.4cm}
    {\large\itshape National Institute of Technology, Tiruchirappalli\par}
    \vspace{1.2cm}

    {\normalsize
    \begin{tabular}{ll}
        \toprule
        \textbf{Name} & \textbf{Roll Number} \\
        \midrule
        Shagnik Sarkar  & 110123102 \\
        Revant Sinha    & 110123088 \\
        Satrajeet Paul  & 110123100 \\
        Mayank Das      & 110123064 \\
        \bottomrule
    \end{tabular}
    \par}
    \vspace{1cm}
    {\normalsize
        \textbf{Dataset:} MIT-BIH Arrhythmia Database (PhysioNet)\\[0.3em]
        \textbf{Framework:} PyTorch 2.x \quad GPU: NVIDIA T4 (Google Colab)\\[0.3em]
        \textbf{Date:} April 2026
    \par}
    \vspace{2cm}

    \begin{abstract}
        \noindent
        This report presents an end-to-end deep learning pipeline for binary
        classification of cardiac arrhythmias from ECG signals.
        Individual heartbeats are extracted from the MIT-BIH Arrhythmia Database
        and transformed into two-dimensional time-frequency scalograms using the
        Continuous Wavelet Transform (CWT) with the Morlet wavelet.  Two
        convolutional neural network (CNN) architectures are trained and evaluated
        on the resulting image representation: (1)~\textbf{SmallNet}, a custom
        lightweight 21-layer CNN with only 3,754 trainable parameters, and
        (2)~\textbf{GoogLeNet}, a pre-trained Inception-v1 network fine-tuned
        via transfer learning with 5,601,954 parameters.  On a balanced
        test set of 3,178 heartbeats, SmallNet achieves an AUC of 82.41\%
        after 12 epochs of CPU training, while GoogLeNet achieves a test
        accuracy of \textbf{95.81\%}, sensitivity of \textbf{96.73\%}, and an
        AUC of \textbf{98.92\%} after 25 GPU epochs.  The results demonstrate
        that the CWT-scalogram approach bridges the gap between 1D biomedical
        signal processing and powerful 2D image classification networks,
        yielding clinically relevant detection performance.
    \end{abstract}

    \vfill
    \rule{\textwidth}{0.5pt}
    {\small All source code is available at:
    \url{https://github.com/Revant-S/BMI_Project}}
\end{titlepage}

\tableofcontents
\newpage

% ====================================================================
\section{Introduction}
% ====================================================================

Cardiovascular disease (CVD) is the leading cause of mortality worldwide,
accounting for approximately 17.9~million deaths annually according to the
World Health Organisation~\cite{who2021cvd}.  Cardiac arrhythmias---irregular
electrical patterns in the heart---are a major contributor to sudden cardiac
arrest and stroke, yet are often asymptomatic and missed in routine check-ups.
The Electrocardiogram (ECG) is the gold-standard non-invasive tool for detecting
arrhythmias; however, manual interpretation of continuous ECG recordings is
time-consuming, requires expert cardiologist skill, and is prone to inter-rater
variability.

Automated arrhythmia detection using machine learning has therefore attracted
significant academic and clinical interest.  Traditional approaches rely on
hand-crafted morphological and statistical features extracted from the R-peak,
QRS complex, and T-wave segments~\cite{de2004automatic}.  Deep learning, in
contrast, enables end-to-end feature learning directly from raw or
lightly-processed signals.

A particularly powerful paradigm is to convert the 1D ECG signal into a 2D
time-frequency image using a wavelet transform, then apply a convolutional
neural network (CNN) designed for image classification~\cite{wang2019ecg}.
This approach combines the rich temporal and spectral information captured by
wavelets with the mature toolset of 2D CNNs---including readily available
pre-trained networks trained on millions of images (transfer learning).

\subsection{Objectives}

\begin{enumerate}[leftmargin=*]
    \item Implement a complete, reproducible pipeline that goes from raw
          MIT-BIH recordings to trained classifiers and evaluation plots.
    \item Convert 1D heartbeat windows into 2D Morlet scalograms and study
          their discriminative quality.
    \item Design and evaluate \textbf{SmallNet}---a custom, resource-efficient
          CNN---as a lightweight clinical-deployment candidate.
    \item Fine-tune \textbf{GoogLeNet} (Inception-v1) via transfer learning as
          a high-accuracy upper-bound reference.
    \item Quantitatively compare both models using clinical metrics:
          Sensitivity, Specificity, Precision, F1-Score, and AUC.
\end{enumerate}

% ====================================================================
\section{Background and Theory}
% ====================================================================

\subsection{The ECG Signal}

The electrocardiogram records the electrical activity of the heart over time.
A single cardiac cycle produces the characteristic PQRST waveform:
\begin{itemize}[leftmargin=*]
    \item \textbf{P-wave}: Atrial depolarisation.
    \item \textbf{QRS complex}: Rapid ventricular depolarisation (the R-peak is
          the sharp positive deflection).
    \item \textbf{T-wave}: Ventricular repolarisation.
\end{itemize}
Arrhythmias manifest as deviations in the timing, shape, or amplitude of these
components.  The MIT-BIH database records signals at $f_s = 360$~Hz with 11-bit
resolution over a $\pm 5$~mV range~\cite{moody2001mitbih}.

\subsection{Beat Segmentation and AAMI Standard}

Individual heartbeats are extracted by locating the R-peak (using ground-truth
annotations provided by cardiologists) and taking a fixed window of
$90 + 160 = 250$ samples ($\approx 694$~ms) centred slightly left of the peak
to capture the pre-excitation onset and the full repolarisation tail.

The Association for the Advancement of Medical Instrumentation (AAMI) EC57
standard \cite{aami2012} defines a recommended label taxonomy for arrhythmia
detectors:
\begin{description}[leftmargin=3em, labelwidth=1.5em]
    \item[N] Normal and bundle-branch block beats
          (\texttt{N, L, R, e, j}).
    \item[S] Supraventricular ectopic beats
          (\texttt{A, a, J, S}).
    \item[V] Ventricular ectopic beats
          (\texttt{V, E}).
    \item[F] Fusion beats (\texttt{F}).
    \item[Q] Unclassifiable or paced beats (excluded from training).
\end{description}
We follow the standard and merge classes S, V, and F into a single
\emph{Abnormal} label, resulting in a binary classification task: Normal vs.\
Abnormal.  Records 102, 104, 107, and 217 are excluded because they contain
paced beats that carry pacemaker artifact morphology unrepresentative of natural
rhythms.

\subsection{Class Imbalance}

The MIT-BIH database is heavily imbalanced: $\approx 90{,}000$ normal beats
versus $\approx 10{,}500$ abnormal beats (8.5:1 ratio).  Training a classifier
on unbalanced data biases it toward predicting the majority class.  We apply
\textbf{random undersampling} of the majority (Normal) class to obtain a
perfectly balanced dataset of $10{,}591 + 10{,}591 = 21{,}182$ beats.

\subsection{Continuous Wavelet Transform (CWT)}

The Continuous Wavelet Transform decomposes a signal $x(t)$ into a
two-dimensional time-scale (time-frequency) representation:
\begin{equation}
    W(a, b) = \frac{1}{\sqrt{|a|}} \int_{-\infty}^{\infty}
              x(t)\,\psi^*\!\left(\frac{t - b}{a}\right) dt
    \label{eq:cwt}
\end{equation}
where $\psi$ is the \emph{mother wavelet}, $a > 0$ is the scale parameter
(inversely proportional to frequency), and $b$ is the translation (time offset).
The resulting matrix $W(a, b)$ is called a \textbf{scalogram}.

The \textbf{Morlet wavelet} (complex sinusoid modulated by a Gaussian) is
chosen because its shape closely resembles the QRS complex, providing optimal
time-frequency resolution for cardiac signals~\cite{addison2005wavelet}:
\begin{equation}
    \psi(t) = \pi^{-1/4} e^{i\omega_0 t} e^{-t^2/2}
    \label{eq:morlet}
\end{equation}
We compute scales $a \in [1, 127]$ (127 frequency bands), producing a
$127 \times 250$ matrix per beat.  This matrix is colour-mapped and resized to
$224 \times 224$ pixels to match the input dimensions expected by standard CNN
architectures.

\subsection{Convolutional Neural Networks}

A CNN applies a series of learnable convolutional filters to extract spatial
feature hierarchies from an input image.  For an input volume $\mathbf{X}$ and
filter bank $\mathbf{W}$, the output of a convolutional layer with bias $b$ and
activation $\sigma$ is:
\begin{equation}
    \mathbf{Y}_{i,j,k} = \sigma\!\left(
        \sum_{m}\sum_{p}\sum_{q} W_{m,p,q,k} \cdot X_{i+p-1,\, j+q-1,\, m}
        + b_k \right)
    \label{eq:conv}
\end{equation}
Batch Normalisation~\cite{ioffe2015batchnorm} and ReLU activations are applied
after each convolutional layer to stabilise training and introduce
non-linearity.  Max-Pooling downsamples the feature maps to progressively
increase the effective receptive field.

\subsection{Transfer Learning}

Transfer learning reuses weight representations learned on a large source task
(ImageNet, $\approx 1.28$~million images, 1000 classes \cite{deng2009imagenet})
and fine-tunes them on a target domain with limited labelled data
\cite{pan2010transfer}.  The rationale is that low-level filters learned on
natural images (edge detectors, texture analysers) are universally useful, even
for medical scalogram images.  Only the final fully-connected head is replaced
to output two classes (Normal / Abnormal).

% ====================================================================
\section{Dataset}
% ====================================================================

\subsection{MIT-BIH Arrhythmia Database}

The MIT-BIH Arrhythmia Database~\cite{moody2001mitbih} was created in 1980 at
the Massachusetts Institute of Technology and Beth Israel Hospital (now BIDMC).
It is freely available via PhysioNet~\cite{goldberger2000physiobank} and
remains the most widely replicated benchmark for arrhythmia detection research,
having been cited in more than 5,000 peer-reviewed publications.

\begin{table}[H]
\centering
\caption{MIT-BIH Arrhythmia Database summary.}
\label{tab:mitbih}
\begin{tabular}{ll}
    \toprule
    \textbf{Property} & \textbf{Value} \\
    \midrule
    Number of half-hour recordings & 48 \\
    Subjects & 47 (25 male, 22 female; age 23--89) \\
    Sampling frequency & 360 Hz \\
    Signal duration per record & $\approx$30 minutes (650,000 samples) \\
    Number of ECG leads & 2 (MLII primary, V5 secondary) \\
    ADC resolution & 11-bit over $\pm 5$ mV \\
    Total annotated beats & $>$110,000 \\
    Annotation method & Two independent cardiologists \\
    \bottomrule
\end{tabular}
\end{table}

\subsection{Annotation Distribution}

Table~\ref{tab:annotations} shows the beat-level annotation counts across the
44 records used in this study (4 paced-beat records excluded).

\begin{table}[H]
\centering
\caption{Beat annotation counts across 44 MIT-BIH records.}
\label{tab:annotations}
\begin{tabular}{clrr}
    \toprule
    \textbf{AAMI Class} & \textbf{Symbol} & \textbf{Description} & \textbf{Count} \\
    \midrule
    \multirow{5}{*}{Normal (N)}
        & N  & Normal sinus beat         & 74,546 \\
        & L  & Left bundle branch block  & 8,075 \\
        & R  & Right bundle branch block & 7,259 \\
        & j  & Nodal escape beat         & 229 \\
        & e  & Atrial escape beat        & 16 \\
    \midrule
    \multirow{6}{*}{Abnormal (S/V/F)}
        & V  & Premature ventricular contraction & 6,903 \\
        & A  & Atrial premature beat             & 2,546 \\
        & F  & Fusion beat                       & 803 \\
        & a  & Aberrated atrial premature beat   & 150 \\
        & E  & Ventricular escape beat           & 106 \\
        & J, S & Nodal / supraventricular premature & 85 \\
    \midrule
    \multicolumn{3}{l}{\textbf{Total Normal (pre-balance)}}  & \textbf{90,125} \\
    \multicolumn{3}{l}{\textbf{Total Abnormal (pre-balance)}}& \textbf{10,593} \\
    \multicolumn{3}{l}{\textbf{After undersampling (1:1)}}   & \textbf{21,182} \\
    \bottomrule
\end{tabular}
\end{table}

% ====================================================================
\section{Methodology}
% ====================================================================

\subsection{Pipeline Overview}

The complete processing pipeline consists of the following sequential stages:

\begin{enumerate}[leftmargin=*]
    \item \textbf{Data Download} (\texttt{download\_data.py}): Fetch 48 MIT-BIH records from PhysioNet via WFDB.
    \item \textbf{Beat Extraction} (\texttt{extract\_beats.py}): Segment individual heartbeats around R-peaks, apply AAMI labelling and class balancing.
    \item \textbf{CWT Scalograms} (\texttt{build\_scalograms.py}): Convert each 1D beat into a $224 \times 224$ RGB scalogram image using the Morlet wavelet.
    \item \textbf{Dataset Loading} (\texttt{dataset.py}): Stratified 70/15/15 split with on-the-fly augmentation.
    \item \textbf{Model Training} (\texttt{train.py}): SGD optimisation with early stopping and checkpoint saving.
    \item \textbf{Evaluation} (\texttt{generate\_results.py}): Confusion matrix, ROC curve, and full metric suite.
\end{enumerate}

\subsection{Beat Extraction}

The following pseudocode summarises the beat extraction procedure:

\begin{lstlisting}[caption={Beat extraction from MIT-BIH records.},label=alg:beat]
for record_id in ALL_RECORDS:
    record = wfdb.rdrecord(record_id)
    ann    = wfdb.rdann(record_id, 'atr')
    signal = record.p_signal[:, 0]  # MLII lead

    for sample, symbol in zip(ann.sample, ann.symbol):
        if symbol in NORMAL_LABELS:
            label = 0
        elif symbol in ABNORMAL_LABELS:
            label = 1
        else:
            continue  # skip non-beat annotations

        start = sample - 90   # 90 samples before R-peak
        end   = sample + 160  # 160 samples after R-peak
        if start < 0 or end > len(signal):
            continue
        beats.append(signal[start:end])  # 250 samples
        labels.append(label)

# Undersample majority class for 1:1 balance
beats, labels = balance_classes(beats, labels)
\end{lstlisting}

\subsection{CWT Scalogram Generation}

For each 250-sample beat vector, the PyWavelets library computes the CWT:
\begin{equation}
    \text{scales} = [1, 2, \ldots, 127], \quad \psi = \text{Morlet}
\end{equation}
The resulting $127 \times 250$ coefficient matrix is logarithmically normalised,
colour-mapped using matplotlib's \texttt{jet} palette, and resized to
$224 \times 224 \times 3$ RGB pixels.  Scalograms are pre-saved to disk as PNG
files ($\approx 24$~KB each, $\approx 500$~MB total) to decouple the
computationally expensive wavelet step from training data loading.

\subsection{Dataset Splits and Augmentation}

A stratified random split (seed $= 42$) partitions the 21,182 scalograms
into:
\[
    \text{Train}: 14{,}827 \;(70\%) \quad
    \text{Validation}: 3{,}177 \;(15\%) \quad
    \text{Test}: 3{,}178 \;(15\%)
\]
Training images receive online augmentation:
\begin{itemize}[leftmargin=*]
    \item Random vertical flip ($p = 0.5$).
    \item Random affine scaling ($\times 0.9$--$\times 1.1$).
    \item Colour jitter (brightness/contrast $\pm 20\%$).
\end{itemize}
All images are normalised to ImageNet statistics
($\mu = [0.485, 0.456, 0.406]$, $\sigma = [0.229, 0.224, 0.225]$) so that
GoogLeNet's pre-trained filters operate in their expected input domain.

\subsection{SmallNet Architecture}

SmallNet is an intentionally compact CNN inspired by the principle that a
well-regularised, domain-specific model can outperform heavyweight architectures
when data volume is limited or inference speed is a clinical priority.

\begin{table}[H]
\centering
\caption{SmallNet layer-by-layer architecture.}
\label{tab:smallnet}
\begin{tabular}{llll}
    \toprule
    \textbf{Block} & \textbf{Operation} & \textbf{Output size} & \textbf{Filters} \\
    \midrule
    Block 1 & Conv $3{\times}3$ -- BN -- ReLU -- MaxPool & $112{\times}112$ & 16 \\
    Block 2 & Conv $3{\times}3$ -- BN -- ReLU -- MaxPool & $56{\times}56$   & 4 \\
    Block 3 & Conv $3{\times}3$ -- BN -- ReLU -- MaxPool & $28{\times}28$   & 8 \\
    Block 4 & Conv $3{\times}3$ -- BN -- ReLU -- MaxPool & $14{\times}14$   & 16 \\
    Head    & AdaptiveAvgPool $\rightarrow$ Flatten       & 16               & --- \\
            & Linear -- ReLU -- Dropout(0.5)              & 64               & 64 \\
            & Linear (output)                             & 2                & 2 \\
    \midrule
    \multicolumn{3}{l}{\textbf{Total trainable parameters}} & \textbf{3,754} \\
    \bottomrule
\end{tabular}
\end{table}

\subsection{GoogLeNet (Inception-v1) Architecture}

GoogLeNet~\cite{szegedy2015goingdeeper} was proposed by Szegedy et al.\ in
2014 and won the ILSVRC 2014 image classification challenge.  Key design
elements include:
\begin{itemize}[leftmargin=*]
    \item \textbf{Inception modules}: Parallel branches of $1{\times}1$,
          $3{\times}3$, $5{\times}5$ convolutions and max-pooling, concatenated
          channel-wise to capture multi-scale features.
    \item \textbf{Auxiliary classifiers}: Two intermediate classification heads
          (AuxClassifier 1 and 2) that inject gradient signals deep into the
          network, reducing vanishing-gradient problems during training.
    \item \textbf{Global Average Pooling}: Replaces the large fully-connected
          layers used in AlexNet/VGG, drastically reducing parameter count.
\end{itemize}

For our binary classification task, the final fully-connected layer is replaced:
\[
    \text{FC}(1024 \rightarrow 1000) \;\longrightarrow\;
    \text{FC}(1024 \rightarrow 2)
\]
The auxiliary heads are similarly adapted ($\text{FC}(\cdot \rightarrow 2)$).
During training, the total loss is:
\begin{equation}
    \mathcal{L} = \mathcal{L}_{\text{main}} +
                  0.3\,\mathcal{L}_{\text{aux1}} +
                  0.3\,\mathcal{L}_{\text{aux2}}
    \label{eq:googleloss}
\end{equation}

\begin{table}[H]
\centering
\caption{GoogLeNet key properties.}
\label{tab:googlenet}
\begin{tabular}{ll}
    \toprule
    \textbf{Property} & \textbf{Value} \\
    \midrule
    Base architecture & Inception-v1 (GoogLeNet) \\
    Pre-training & ImageNet (1.28M images, 1000 classes) \\
    Inception modules & 9 \\
    Auxiliary classifiers & 2 \\
    Total depth & 22 layers \\
    Total trainable parameters & 5,601,954 \\
    Input size & $224{\times}224{\times}3$ \\
    Output classes & 2 (Normal / Abnormal) \\
    \bottomrule
\end{tabular}
\end{table}

\subsection{Training Configuration}

Both models are trained with identical optimiser settings to ensure a fair
comparison.

\begin{table}[H]
\centering
\caption{Shared training hyperparameters.}
\label{tab:hyperparams}
\begin{tabular}{ll}
    \toprule
    \textbf{Hyperparameter} & \textbf{Value} \\
    \midrule
    Optimiser         & SGD with momentum \\
    Learning rate     & $\eta = 1 \times 10^{-4}$ \\
    Momentum          & 0.9 \\
    Loss function     & Cross-Entropy \\
    Batch size        & 32 \\
    Max epochs        & 50 (with early stopping) \\
    Early stop patience & 10 epochs \\
    Random seed       & 42 \\
    Hardware          & NVIDIA T4 GPU (Google Colab, 15 GB VRAM) \\
    \bottomrule
\end{tabular}
\end{table}

% ====================================================================
\section{Experiments and Results}
% ====================================================================

\subsection{SmallNet Training}

SmallNet was trained for 12 epochs; the best validation checkpoint was saved at
epoch 2 (val\_acc $= 76.14\%$).  After epoch 2 the validation accuracy
plateaued between 70--76\%, indicating that the model's capacity (3,754
parameters) was saturated---the 4-block feature extractor exhausted its ability
to discriminate the complex multi-scale wavelet patterns.

\begin{table}[H]
\centering
\caption{SmallNet training progression (selected epochs).}
\label{tab:smallnet_hist}
\begin{tabular}{crrrr}
    \toprule
    \textbf{Epoch} & \textbf{Train Loss} & \textbf{Train Acc (\%)} &
    \textbf{Val Loss} & \textbf{Val Acc (\%)} \\
    \midrule
    1  & 0.7015 & 50.53 & 0.6884 & 49.98 \\
    2  & 0.6954 & 51.50 & 0.6850 & \textbf{76.14} \\
    4  & 0.6867 & 54.31 & 0.6745 & 70.16 \\
    6  & 0.6716 & 60.79 & 0.6588 & 70.07 \\
    8  & 0.6552 & 64.92 & 0.6393 & 73.15 \\
    10 & 0.6427 & 66.65 & 0.6208 & 72.77 \\
    12 & 0.6302 & 67.92 & 0.6020 & 71.92 \\
    \bottomrule
\end{tabular}
\end{table}

\subsection{GoogLeNet Training}

GoogLeNet converged smoothly over 25 epochs.  Transfer learning from ImageNet
provided a powerful feature initialisation: validation accuracy crossed 90\%
within 5 epochs and reached 95.85\% by epoch 25.

\begin{table}[H]
\centering
\caption{GoogLeNet training progression (selected epochs).}
\label{tab:googlenet_hist}
\begin{tabular}{crrrr}
    \toprule
    \textbf{Epoch} & \textbf{Train Loss} & \textbf{Train Acc (\%)} &
    \textbf{Val Loss} & \textbf{Val Acc (\%)} \\
    \midrule
    1  & 0.5916 & 70.43 & 0.4799 & 81.11 \\
    5  & 0.3148 & 86.83 & 0.2555 & 90.24 \\
    10 & 0.2257 & 91.45 & 0.1741 & 93.26 \\
    15 & 0.1848 & 93.07 & 0.1441 & 94.52 \\
    20 & 0.1560 & 94.25 & 0.1247 & 95.56 \\
    25 & 0.1387 & 94.94 & 0.1183 & \textbf{95.85} \\
    \bottomrule
\end{tabular}
\end{table}

\subsection{Test Set Performance}

Table~\ref{tab:results} presents the final evaluation on the held-out test set
(3,178 samples, equal class balance).

\begin{table}[H]
\centering
\caption{Test set performance comparison.}
\label{tab:results}
\begin{tabular}{lcc}
    \toprule
    \textbf{Metric} & \textbf{SmallNet} & \textbf{GoogLeNet} \\
    \midrule
    Accuracy (\%)    & 75.77 & \textbf{95.81} \\
    Sensitivity (\%) & 79.36 & \textbf{96.73} \\
    Specificity (\%) & 72.18 & \textbf{94.90} \\
    Precision (\%)   & 74.05 & \textbf{94.99} \\
    F1-Score (\%)    & 76.61 & \textbf{95.85} \\
    AUC (\%)         & 82.41 & \textbf{98.92} \\
    \midrule
    True Positives   & 1,261 & 1,537 \\
    True Negatives   & 1,147 & 1,508 \\
    False Positives  & 442   & 81 \\
    False Negatives  & 328   & 52 \\
    \midrule
    Trainable Params & 3,754 & 5,601,954 \\
    Training Epochs  & 12 (CPU) & 25 (GPU) \\
    \bottomrule
\end{tabular}
\end{table}

\subsection{Discussion}

\paragraph{GoogLeNet performance.}
GoogLeNet's 98.92\% AUC demonstrates near-perfect class separation on the
test set.  The high sensitivity (96.73\%) is especially important in medical
screening: out of 1,589 truly abnormal beats in the test set, only 52 were
missed (False Negatives).  The low false-negative rate is critical in clinical
practice where missed arrhythmias can be life-threatening.

\paragraph{SmallNet analysis.}
SmallNet's plateau at $\approx 76\%$ accuracy illustrates the classical
\emph{capacity bottleneck} problem.  With only 3,754 parameters, the network
is fundamentally limited in the complexity of the decision boundary it can
represent.  Nonetheless, its 82.41\% AUC indicates it has learned a
statistically meaningful discrimination between Normal and Abnormal scalogram
patterns---a remarkable result given its minimal computational budget.

\paragraph{Intra-patient vs.\ inter-patient evaluation.}
A key limitation of the current evaluation is the use of an
\emph{intra-patient} split: beats from the same patient may appear in both the
training and test sets.  Since individual heartbeat morphology has a
patient-specific ``fingerprint'', the model can leverage this memorisation to
inflate accuracy.  Published results using an \emph{inter-patient} split
(training on a disjoint set of patients from the test set) typically report
5--10 percentage points lower accuracy.  A rigorous clinical evaluation would
require the inter-patient protocol.

% ====================================================================
\section{Conclusion}
% ====================================================================

This project presents a complete, modular, and reproducible pipeline for binary
ECG arrhythmia classification.  The key contribution is the application of the
Continuous Wavelet Transform to convert 1D heartbeat signals into
two-dimensional Morlet scalogram images, enabling the direct application of
powerful image CNN architectures.

GoogLeNet, fine-tuned via transfer learning, achieves clinically relevant
performance: 95.81\% accuracy, 96.73\% sensitivity, and 98.92\% AUC on the
balanced MIT-BIH test set.  SmallNet, with only 3,754 parameters, provides a
lightweight baseline that can serve as a starting point for resource-constrained
embedded medical devices.

\paragraph{Future work.}
\begin{itemize}[leftmargin=*]
    \item Adopt a strict \emph{inter-patient} train/test split following the
          AAMI DS1/DS2 protocol for more realistic generalisation estimates.
    \item Explore multi-class classification beyond binary Normal/Abnormal.
    \item Investigate model compression techniques (pruning, quantisation) to
          deploy SmallNet on microcontrollers.
    \item Integrate Grad-CAM visualisations to interpret which scalogram
          regions are driving the GoogLeNet predictions.
    \item Extend the dataset with additional PhysioNet databases (e.g.,
          INCART, European ST-T) to improve generalisation.
\end{itemize}

% ====================================================================
\section*{Acknowledgements}
% ====================================================================

\noindent
The MIT-BIH Arrhythmia Database was obtained from PhysioNet
(\url{https://physionet.org/content/mitdb/}).  Model training was performed
on Google Colab's free-tier NVIDIA T4 GPU.  The PyTorch, torchvision,
PyWavelets, wfdb, scikit-learn, and matplotlib open-source libraries were
used throughout.

% ====================================================================
\begin{thebibliography}{99}

\bibitem{who2021cvd}
World Health Organization (2021).
\textit{Cardiovascular diseases (CVDs) fact sheet}.
\url{https://www.who.int/news-room/fact-sheets/detail/cardiovascular-diseases}

\bibitem{de2004automatic}
de~Chazal, P., O'Dwyer, M., and Reilly, R.~B. (2004).
Automatic classification of heartbeats using ECG morphology and heartbeat
interval features.
\textit{IEEE Transactions on Biomedical Engineering}, 51(7):1196--1206.

\bibitem{wang2019ecg}
Wang, T., Lu, C., Sun, Y., Yang, M., Liu, C., and Ou, C. (2019).
Automatic ECG classification using continuous wavelet transform and
convolutional neural network.
\textit{Entropy}, 23(1):119.

\bibitem{moody2001mitbih}
Moody, G.~B. and Mark, R.~G. (2001).
The impact of the MIT-BIH arrhythmia database.
\textit{IEEE Engineering in Medicine and Biology Magazine}, 20(3):45--50.

\bibitem{goldberger2000physiobank}
Goldberger, A.~L. et al.\ (2000).
PhysioBank, PhysioToolkit, and PhysioNet: Components of a new research resource
for complex physiologic signals.
\textit{Circulation}, 101(23):e215--e220.

\bibitem{aami2012}
ANSI/AAMI EC57:2012.
\textit{Testing and reporting performance results of cardiac rhythm and ST
segment measurement algorithms}.
Association for the Advancement of Medical Instrumentation, Arlington, VA.

\bibitem{addison2005wavelet}
Addison, P.~S. (2005).
Wavelet transforms and the ECG: A review.
\textit{Physiological Measurement}, 26(5):R155--R199.

\bibitem{ioffe2015batchnorm}
Ioffe, S. and Szegedy, C. (2015).
Batch normalization: Accelerating deep network training by reducing internal
covariate shift.
\textit{Proceedings of ICML}, 37:448--456.

\bibitem{deng2009imagenet}
Deng, J. et al.\ (2009).
ImageNet: A large-scale hierarchical image database.
\textit{IEEE CVPR}, 248--255.

\bibitem{pan2010transfer}
Pan, S. J. and Yang, Q. (2010).
A survey on transfer learning.
\textit{IEEE Transactions on Knowledge and Data Engineering},
22(10):1345--1359.

\bibitem{szegedy2015goingdeeper}
Szegedy, C. et al.\ (2015).
Going deeper with convolutions.
\textit{IEEE CVPR}, 1--9.

\end{thebibliography}

\end{document}
